\chapter*{Preface}
\addcontentsline{toc}{chapter}{Preface}
\markboth{Preface}{Preface}

Combinatorics and graph theory are two subjects that keep
getting bundled together because they keep needing each other: counting
carefully, and the structures --- graphs --- that show up whenever
counting stops being about isolated objects and starts being about how
objects connect. This booklet covers counting principles, generating
functions, inclusion-exclusion, the pigeonhole principle and Ramsey
theory, then graphs, trees, connectivity, coloring, planarity, and
network flows.

It depends on Booklets~01 and~02 for its proof techniques and its
language of sets and relations, and not much else --- which is part of
why it sits here, early, rather than after the analysis booklets.
